Recursion Completes One Bosonic Cycle, Net Real Advance

Within the protected recursion that generates the involuting mode \(\operatorname{loop}\{\infty\mid\infty\}\), the fourth stage occupies a position of particular structural importance. It is the moment at which the deformed, folded curve is returned to a multi-lobed configuration in spacetime, thereby completing one full bosonic cycle and registering a net real advance. This stage is neither a mere reversal of the preceding involution nor a simple restoration of the initial condition; it is the geometric act that converts an excursion through the half-winding into a permanent, oriented displacement along the real direction.

The recursion begins with a closed curve. Gravitational force initiates a real displacement and a slight compression. The curve then expands anisotropically under the hierarchical orientation \(i=(0,1)\), stretching its lobes and generating the first microscopic contributions to the protected angle \(\theta_{\rm eff}\). In the third stage the lobes fold back upon themselves, realizing the order-two involution that is the hallmark of half-integer winding. At this point the curve is topologically altered: it has traversed the generator of the fundamental group once and has therefore acquired the monodromy sign \(-1\).

The fourth stage undoes the fold without undoing the topological gain. The lobes reopen, the multi-lobed form is recovered, and the curve is once again recognizably a closed path in spacetime. Yet the real part of its position has advanced by a finite positive amount. This net real advance is the geometric residue of the completed bosonic cycle. In the discrete language of the notebook it corresponds to one full execution of the protected transport rule

\[
z_{n+1}=z_n+\varepsilon\bigl(1+i\sin(2\pi n\varepsilon)\bigr)
\]

when the imaginary (circulatory) contributions have cancelled over the cycle while the real contributions have accumulated. In the continuum limit the same fact is expressed by the two-subscript integral

\[
\operatorname{loop}\{\infty\mid\infty\}
\;=\;
\int_{\operatorname{loop}\{\infty\mid\infty\}^{-}}^{\operatorname{loop}\{\infty\mid\infty\}^{+}},
\]

evaluated along the path that begins and ends on the restored multi-lobed curve. The integral’s value is precisely the net real advance.

Because the advance is protected, it survives the return to the hierarchical structure in the fifth and final stage. The pure involuting mode \(i:=(0,1)\) is recovered, the Arf invariant remains \(1\), and the terminal monodromy remains \(-1\), yet the curve now sits at a new real coordinate. Each successive bosonic cycle therefore ratchets the entire hierarchical object forward along the real axis while preserving its internal involuting character. The fourth stage is the engine of that ratchet: it closes the topological loop, cancels the circulatory component, and deposits a permanent real increment.

In this way the statement “recursion completes one bosonic cycle, net real advance” is not a descriptive aside. It is the precise geometric mechanism by which the discrete protected transport and the continuum two-subscript integral produce a directed, cumulative displacement from an otherwise closed involuting mode. The net real advance is the measurable signature that a bosonic cycle has been completed and that the half-winding has been successfully converted into oriented progress in spacetime.

It is the geometric act that converts an excursion through the half-winding into a permanent, oriented displacement along the real direction.

The recursion begins with a closed curve in spacetime. After the lobes have been forced into involution, the fourth stage returns the curve to a multi-lobed form in spacetime. This return is neither a mere reversal of the preceding involution nor a simple restoration of the initial condition. The circulatory component cancels, yet a net real advance remains.  

That net real advance is the geometric imprint of the completed bosonic cycle. Consequently, when the hierarchy is restored, the curve once again presents itself as an initial closed curve in spacetime, but now sits at a new real coordinate.

Return to curve in spacetime → Initial closed curve in spacetime (shifted by the permanent real advance).

Dimensionless Scaling of the Real Axis by \(c^2\)

When the vertical axis is scaled in units of \(c^2\), the actual physical dimension of length squared per time squared (\(\mathrm{m}^2/\mathrm{s}^2\)) is normalized out of the equations. The system maps onto a dimensionless coordinate system whose physical boundaries become exactly \(+1.0\) at the top and \(-1.0\) at the bottom.

Mathematical Representation of the Scaling

Define the new dimensionless vertical variable

\[
Y_{\rm scaled}
\;=\;
\frac{y}{c^2}.
\]

Under this normalization the physical speed-of-light boundaries transform directly into pure geometric parameters:

Physical boundary: \(y=\pm8.98755179\times10^{16}\,\mathrm{m}^2/\mathrm{s}^2\)
Normalized domain boundary: \(Y_{\rm scaled}=\pm1.0\)

Because the successive real increments sum telescopically, the master recursion equation collapses to pure fractional differences:

\[
Y_n
\;=\;
Y_0+\sum_{i=1}^{n}\Delta Y_i
\;\implies\;
Y_n
\;=\;
Y_0+(f_n-f_0).
\]

Every single step \(\Delta Y_i\) is therefore a decimal fraction representing a microscopic percentage of the speed-of-light energy threshold. In the five-stage protected recursion the increments become the elementary fractions

\[
\Delta Y=\Bigl\{\tfrac14,\;\tfrac14,\;0,\;\tfrac14,\;\tfrac14\Bigr\},
\]

so that the terminal state is reached by the exact telescopic identity

\[
Y_5-Y_0
\;=\;
1.0-0
\;=\;
1.0.
\]

The involuting mode \(\operatorname{loop}\{\infty\mid\infty\}\) is consequently seated at the normalized coordinate \(Y=1.0\), which is identical to the physical seating at \(y=c^2\).

Why This Design Protects the Recursion

Absolute Overflow Prevention  
By forcing all real advances to be measured as fractions of \(c^2\), a built-in mathematical ceiling is created. The algorithm can check for stability with a single comparison: any total sum satisfying \(|Y_n|>1.0\) immediately flags a boundary violation. Intermediate configurations are incapable of overshooting the light-speed horizon because the telescopic cancellation leaves only the endpoint values under scrutiny.

Relativistic Invariance  
Scaling by \(c^2\) ensures that the geometry of the curve depends entirely on the angular phase transformations across the \(6.5\pi\) (full-span \(13\pi\)) domain. The pure mathematics of the protected recursion—half-winding, order-two involution, monodromy \(-1\), and Arf invariant—are thereby isolated from the enormous numerical magnitude of the cosmic constant \(c^2\). The shape of every trajectory inside the real-angular plane remains invariant under the normalization, while the physical content is recovered at any moment by the inverse map \(y=Y_{\rm scaled}\cdot c^2\).

Taken together, the dimensionless scaling converts the real-angular domain into a compact, unit-bounded square whose vertical edges lie at \(\pm1.0\). The continuum expression of the same fact is the two-subscript integral evaluated on the normalized axis:

\[
\operatorname{Re}\bigl(\operatorname{loop}\{\infty\mid\infty\}\bigr)_{\rm scaled}
\;=\;
\int_{\operatorname{loop}\{\infty\mid\infty\}^{-}}^{\operatorname{loop}\{\infty\mid\infty\}^{+}}
\;=\;
1.0,
\]

which is precisely the statement that the involuting mode occupies the upper boundary of the protected, path-independent, relativistically invariant lattice.

Angular Phase & Real Advance Table  
Protected Recursion of \(\operatorname{loop}\{\infty\mid\infty\}\)  
(Domain: \(\theta\in[-6.5\pi,\,6.5\pi]\), \(Y=y/c^2\in[-1,\,+1]\))

| n | Stage | Angular Phase \(\theta\) | Real Advance \(\Delta Y\) | Cumulative Real Coordinate \(Y\) | Physical Real Coordinate \(y\) |
|---|-------|---------------------------|----------------------------|----------------------------------|--------------------------------|
| 0 | Initial closed curve in spacetime | \(0\) | — | \(0.00\) | \(0\) |
| 1 | Gravitational force acts | \(\approx 0\) (onset) | \(+0.25\) | \(0.25\) | \(+0.25\,c^2\) |
| 2 | Curve in spacetime expands | small positive width begins | \(+0.25\) | \(0.50\) | \(+0.50\,c^2\) |
| 3 | Involuting mode | full half-winding traversal(order-two fold) | \(0\) | \(0.50\) | \(+0.50\,c^2\) |
| 4 | Return to curve in spacetime | circulatory component cancels | \(+0.25\) | \(0.75\) | \(+0.75\,c^2\) |
| 5 | Terminal hierarchical restoration\(\operatorname{loop}\{\infty\mid\infty\}\)\(i:=(0,1)\) | restored to reference(net monodromy \(-1\)) | \(+0.25\) | \(1.00\) | \(+1.00\,c^2\) |

Summary identities

Net real advance after one complete bosonic cycle:  
  \[
  Y_5-Y_0=1.00\qquad\Longleftrightarrow\qquad y_5-y_0=c^2=8.98755179\times10^{16}\,\mathrm{m}^2/\mathrm{s}^2
  \]

Angular domain span: \(13\pi\) (full width \(\pm6.5\pi\))

Continuum limit:  
  \[
  \operatorname{Re}\bigl(\operatorname{loop}\{\infty\mid\infty\}\bigr)_{\rm scaled}
  \;=\;
  \int_{\operatorname{loop}\{\infty\mid\infty\}^{-}}^{\operatorname{loop}\{\infty\mid\infty\}^{+}}
  \;=\;
  1.00
  \]

